潜变量模型想用看不见的 \(Z\) 解释看得见的 \(X\)，写下来很干净，一算就卡住：边缘似然是一个高维积分，后验的归一化常数正是那个积分。VAE 的对策是不去算它——拿一个可训练的分布 \(q_\phi(z\given x)\) 逼近后验，把``最大化似然''换成``最大化一个能求梯度的下界''。

这一步换来的东西，也带来了三笔账，这一章分三篇分别结清。第一笔是\emph{下界的松紧}：ELBO 与真实对数似然之间隔着一个未知的 KL 间隙，所以 ELBO 上升不等于生成模型变好。第二笔是\emph{梯度怎么穿过采样}：目标是对一个依赖参数的分布取期望，梯度不能直接穿过随机节点，得分函数估计与重参数化是两条走法，方差和适用范围都不同。第三笔是\emph{目标允许的退化解}：ELBO 里的 KL 项在惩罚隐变量被使用，当解码器足够强时，把隐变量彻底闲置反而是最优的。

贯穿三篇的判断是同一条：把难算的量换成可优化的替代品，替代品的形状会反过来决定你实际得到什么。

\subsection{怎么读这一章}

第一篇是入口，生成机制、不可积后验和 ELBO 分解都在里面，跳过它后两篇无从谈起。第二篇是梯度估计的技术核心，如果只想先跑通一个 VAE，可以在读完重参数化的公式后暂停，把方差比较和离散松弛留到需要处理离散隐变量时再回来。第三篇在真正训练出一个``损失很漂亮但隐空间是空的''模型之后读，体会最深。

三条不能弄混的界线：ELBO 不是似然，重构误差不是对数似然，KL 不为零不等于隐变量真的有用。

\subsection{本章篇目}

以下篇目按概念依赖排列；若从具体问题进入，也可以先打开最接近当前困惑的一篇，再沿篇内链接回补。

\begin{itemize}
\tightlist
\item \hyperref[sec:14-Deep-Learning/06-Variational-Autoencoders/01]{``VAE 把不可积后验变成可训练下界''}：为什么边缘似然和后验被同一个积分堵死，Jensen 不等式如何给出 ELBO，重构项与 KL 项为什么是同一个分解的两半。
\item \hyperref[sec:14-Deep-Learning/06-Variational-Autoencoders/02]{``重参数化让随机节点可以反向传播''}：参数长在期望下标上时梯度为何断掉，得分函数估计与重参数化各自要赔上什么，交换求导与积分需要哪些条件，以及离散隐变量为什么只能接受有偏的连续松弛。
\item \hyperref[sec:14-Deep-Learning/06-Variational-Autoencoders/03]{``后验坍缩与表示能力的失衡''}：坍缩为什么是 ELBO 的结构性后果而非训练故障，KL 项作为信息量上界给出的率失真读法，四种缓解手段分别在改路径、改目标、改约束还是改模型。
\end{itemize}

读完这一章，VAE 的位置就定下来了：它选了一个特定的散度、一个特定的下界、一族特定的近似后验，三个选择共同决定了它擅长什么、在哪里失败。下一个单元把 VAE 与生成对抗网络、流模型、扩散模型放进同一个框架，逐个问它们实际在最小化哪个散度——那个答案能解释为什么 VAE 的样本偏糊而对抗训练的样本偏尖锐却容易丢模式。
